% !TeX root = ../../thesis.tex
\chapter{Acknowledgements}\label{ch:acknowledgements}

\instructionsintroduction

% Illustration on how to refer to your papers when using biblatex
% (see second line in thesis.tex to activate biblatex)
%\definecolor{shadecolor}{gray}{0.90}
%\begin{shaded}
%A nice citation
%\end{shaded}
%This chapter was previously published as:\\
%\fullcite{VandenBroeck2011IJCAI}
%\newpage


% Some dummy code to make sure bibtex does not complain.
%Illustration of how to include citations \cite{Meert2011PhD} and \cite{VandenBroeck2011IJCAI}.
%Lorem ipsum dolor sit amet, consetetur sadipscing elitr, sed diam nonumy eirmod tempor invidunt ut labore et %dolore magna aliquyam erat, sed diam voluptua. At vero eos et accusam et justo duo dolores et ea rebum. Stet %clita kasd gubergren, no sea takimata sanctus est Lorem ipsum dolor sit amet.

%And yet another citation \cite{FrRo2010Diffusion}.

\markboth{\uppercase{Acknowledgements}}{\uppercase{Acknowledgements}}
%\markright{\uppercase{Scoring metrics of codon usage and of tRNA adaptation}}

% Some dummy code to get at least 1 entry in the list of
% abbreviations.
%\newglossaryentry{md}{name={MD},description={molecular dynamics}}
%Introducing an acronym: \gls{md}.
%{\small
Before starting my scientific career, I was involved in corporate activities, working for a variety of small and medium enterprises. My curiosity and appetite for learning drove me to enroll in multiple STEM academic programs in parallel, all while working full-time as an employee. Students usually engage in \textit{extra-curricular} activities—I, on the other hand, was a corporate employee engaging in \textit{extra-corporate-lar} activities, expanding both my skill set and my overtime hours. One might say I had a \textit{double major} in business and curiosity.

During this long corporate journey, I had the chance to meet highly skilled professionals, committed employees, and even business angels. I was never bored and always enjoyed looking up at the sky—while keeping my feet firmly on the ground or, quite literally, in the mud, as I worked with wastewater treatment sludge and landfill leachates.

Many former corporate colleagues and bosses have described me as someone who blends a highly theoretical and conceptual education with practical, hands-on expertise. At my core, I find true fulfillment in learning and satisfying my insatiable curiosity for new knowledge.

After corporate workers and business angels, I met my \textit{research angel}: Professor Liesbet Geris. Lies, I wish to express my deepest gratitude for being the first person to take my unconventional academic path seriously and for having the courage to invest in me—without bias, despite my non-traditional trajectory as a senior candidate, with no scientific seniority, in his fifties. I sincerely appreciate your trust and, most importantly, your generous research funding, which has supported me for seven very happy years. I have told you this before, but I say it again: you are a modern Renaissance patron, championing both science and the arts. Your willingness to take a risk on an interdisciplinary PhD project outside your primary research track demonstrates your vision as a principal investigator. At least we shared a common foundation in modeling in computational biology. I am also deeply appreciative of the guidance, rigor, and scientific integrity you have instilled in me through our discussions—whether in your office or virtually, often extending into evenings, weekends, and even late nights on \textit{Slack}. Your expertise in modeling and digital twin development has been invaluable. You were among the pioneers in the field, shaping how researchers conceptualize \textit{in silico} representations of biological systems with increasing complexity. I also recognize and value the effort you have dedicated to teaching me how to write more concise papers. I promise to keep improving in that area! Your strategic insights on selecting editorial boards and navigating peer review have been incredibly instructive. I am proud and grateful to have worked with you and to have witnessed the numerous prestigious grants, awards and \textit{honoris causa} distinctions you have received across Europe.

Selecting the right PhD research topic is of utmost importance, as it will define several years of work. Professor Michel Georges, former director of GIGA, once advised PhD candidates in an introductory lecture at the doctoral school that the most effective way to begin a doctoral project is to identify the boundary between what is known and what remains unknown. Once this frontier is found, the task of the researcher is to move to the edge and push it further.

How did I find the topic and how did I move to the edge ?

In November 2018, during a discussion at the GIGA doctoral school, Professor Pierre Close introduced me to one of his research interests: enzymatic modifications of nucleotide 34 in tRNAs and their effects on protein elongation rates. He explained how transcripts with different codon usage are translated at varying efficiencies, leading to shifts in the proteome landscape. In melanoma and lung cancer contexts, these modifications to U34 or A34-tRNAs alter the oncoproteome. A seemingly small chemical modification in the anticodon loop of a tRNA molecule could influence the translation of the entire proteome.

At first glance, the underlying molecular mechanisms appeared bafflingly complex—at the very least, fascinating. This topic immediately captivated me, as it challenged my background in chemical and bioengineering. Pierre, thank you for bringing me on this journey. After several weeks of literature review, I became convinced that the key to advancing knowledge in this field lay in addressing its complexity through computational modeling—systematically integrating the multiple contributing factors described in the literature. Given that at least five interconnected variables were involved, tackling this question required a truly interdisciplinary and integrative approach. I also want to express my gratitude to Pierre Close coworkers: Dr. Francesca Rapino, Dr. Arnaud Blomme and Marine Leclercq for sharing their datasets.

I would like to express my retrospective gratitude to my high school physics teacher, Madame Marie-Thérèse Henrotin, whose curiosity and passion for teaching left a lasting impression on me. A special thank you goes to my chemistry teacher, Monsieur Alain Capouillez, an autodidact who introduced me to Henry Eyring’s transition state theory of catalysis when I was only 16. This was not part of the standard curriculum, and I am immensely grateful that he chose to go beyond it. He would have certainly appreciated the mechanochemistry embedded in the transition-state theory applied in this thesis to study the ribozyme-catalyzed kinetics of peptide bond formation.

I would also like to extend my sincere thanks to the members of my thesis jury for their time and interest in my research. In particular, I am grateful to Professor Eveline Lescrinier for her co-endorsement—alongside Professor Liesbet Geris—of my future postdoctoral project. I also extend my appreciation to Professor André Matagne, Dr. Frédéric Kerff and Dr. Arnaud Vanden Broeck. A sincere thank you to professor André Matagne for welcoming me as a member of the Belgian Biophysical Society. I truly appreciate the opportunity to be part of this community and to engage with fellow researchers in the field of biophysics. I look forward to contributing to the society’s activities and benefiting from the exchange of ideas and collaborations.
I warmly thank Professor Gerben Menschaert for the technical advice, valuable discussions, and insightful exchanges we had at Ghent University regarding Ribo-Seq technology and its data analysis. My gratitude also goes to Professor Dick de Ridder for organizing the Bioinformatics Summer School sessions at Wageningen University with his team, as well as for his pioneering work over a decade ago on modeling protein synthesis using TASEP.

Additionally, I would like to acknowledge Professor Tamir Tuller of Tel Aviv University. In these challenging times, researchers around the world face difficulties due to conflicts, the rise of authoritarian regimes or repressive governance beyond their control. I deeply regret the isolation that many scientists endure, ranging from professional exclusion to, in the worst cases, acts of censorship, restrictions on freedoms or even war-related casualties. It is my hope that scientific collaboration and open discourse will continue to transcend political and ideological barriers.

I would like to express my gratitude to all my fellow colleagues, PhD candidates and postdocs in Lies' Biomechanics group with whom I spent joyful moments: Bernard Staumont, Sophie Bekisz, Loic Comeliau, Luiz Ladeira, Alessio Gamba, Margaryta Ivanets, Sophie Nguyen, Ehsan Sadeghian, Fernando Perez Boerema, Bingbing Liang, Satanik Mukherjee, Mohammad Mehrian, Morgane Germain, Niki Loverdou, Raphaelle Lesage, Majid Nazemi, Laura Lafuente Gracia, Mojtaba Barzegari Shankil, Tim Herpelinck, Lisanne Groeneveldt, Sourav Mandal, Oriana de Becker, Ayse Kose, Tom Verbraeken, Mervenaz Sahin, Edoardo Borgiani, Gabriele Nasello, Claire Vilette, Ahmad Alminnawi, Pieter Ansoms, Liesbeth Ory.

I want to thank Isabelle Rausin for being the best mother of our three daughters and for her valuable networking sharing which gave me the opportunity to meet Professor Liesbet Geris.

Finally, and most importantly, I want to express my heartfelt gratitude to my family—my father, and my late mother, whose unwavering enthusiasm and support still inspire me. To Laurence, for her encouragement and patience, and to my three daughters, Emmy, Sarah and Lara, whose achievements make me so proud and with whom I share a deep and joyful complicity, making them my most cherished and passionate companions in life.

This work was supported by the FWO EOS grant no.30480119 (Joint-t-against-Osteoarthritis) in Belgium, the European Research Council under the European Union's Horizon 2020 Framework Program (H2020/2014-2020) ERC grant no.772418 (INSITE) and the ERC grant no.101088919 INStant CARMA In Silico Trials for Cartilage Regenerative Medicine Applications.

Computational resources have been provided by the Consortium des Équipements de Calcul Intensif (CÉCI), funded by the Fonds de la Recherche Scientifique de Belgique (F.R.S.-FNRS) under grant no.2.5020.11 and by the Walloon Region.
%} %end small



%%%%%%%%%%%%%%%%%%%%%%%%%%%%%%%%%%%%%%%%%%%%%%%%%%
% Keep the following \cleardoublepage at the end of this file,
% otherwise \includeonly includes empty pages.
\cleardoublepage

% vim: tw=70 nocindent expandtab foldmethod=marker foldmarker={{{}{,}{}}}
